\abstract{РЕФЕРАТ}

\setlength{\parindent}{1.25cm}

Объем работы равен \formbytotal{lastpage}{страниц}{е}{ам}{ам}. Работа содержит \formbytotal{figurecnt}{иллюстраци}{ю}{и}{й}, \formbytotal{tablecnt}{таблиц}{у}{ы}{}, \arabic{bibcount} библиографических источников и \formbytotal{числоПлакатов}{лист}{}{а}{ов} графического материала. Количество приложений – 2. Графический материал представлен в приложении А. Фрагменты исходного кода представлены в приложении Б.

Перечень ключевых слов: корпоративный мессенджер, система, сервер, web-приложение, группы, каналы, сообщения, администратор, пользователь база данных, пользователи, сообщения, группы, безопасность, шифрование, файлы, логирование, API, производительность, администрирование, API, аутентификация, авторизация, приложение, модуль.

Объектом разработки является корпоративный мессенджер, обеспечивающий безопасное взаимодействие сотрудников внутри организации, защиту паролей, управление доступом и интеграцию с внутренними системами компании.

Целью выпускной квалификационной работы является создание корпоративного коммуникационного инструмента, обеспечивающего обмен сообщениями, файлами и уведомлениями между сотрудниками, а также возможность администрирования групп, настройки прав и контроля доступа.

В процессе разработки серверной части системы была выделена архитектура API, разработаны модули обработки данных, реализована логика управления пользователями, группами и каналами, а также механизм шифрования паролей пользователей.

При разработке серверной части мессенджера использовались Python, Webob, SQLite.

\selectlanguage{english}
\abstract{ABSTRACT}
  
The volume of work is \formbytotal{lastpage}{page}{}{s}{s}. The work contains \formbytotal{figurecnt}{illustration}{}{s}{s}, \formbytotal{tablecnt}{table}{}{s}{s}, \arabic{bibcount} bibliographic sources and \formbytotal{числоПлакатов}{sheet}{}{s}{s} of graphic material. The number of applications is 2. The graphic material is presented in annex A. The layout of the site, including the connection of components, is presented in annex B.

List of keywords: corporate messenger, system, server, web application, groups, channels, messages, administrator, user database, users, messages, groups, security, encryption, files, logging, API, performance, administration, API, authentication, authorization, application, module.

The object of development is a corporate messenger that ensures secure interaction of employees within the organization, password protection, access management and integration with the company's internal systems.

The goal of the final qualification work is to create a corporate communication tool that ensures the exchange of messages, files and notifications between employees, as well as the ability to administer groups, configure rights and access control.

During the development of the server part of the system, the API architecture was highlighted, data processing modules were developed, the logic for managing users, groups and channels was implemented, as well as the mechanism for encrypting user passwords.

Python, Webob, SQLite were used in the development of the server part of the messenger.

\selectlanguage{russian}
